\section*{Chapitre \ref{ch:pl:signal-serving-and-parameter-management} --- Service de signaux et gestion des paramètres}

\begin{solution}{exo:pl:signal-serving-and-parameter-management:1}
50 points de base sur 4~millions de dollars font 20\,000~\$, soit 200 actions à 100~\$~: deux lots. 50~\% font 2~millions de dollars, soit 20\,000 actions~: deux cents lots.
\end{solution}

\begin{solution}{exo:pl:signal-serving-and-parameter-management:2}
La hausse est un pas de 25 points de base (0.0025), au-dessus des 20 qui exigent un second approbateur~; la correction de 75 à 60 est un pas de 15, inférieur
à ce seuil.
\end{solution}

\begin{solution}{exo:pl:signal-serving-and-parameter-management:3}
0.7, version 3, fraîche~: la dernière valeur a été publiée à 60~secondes, soit 340~secondes plus tôt, dans la limite de cinq minutes.
\end{solution}

\begin{solution}{exo:pl:signal-serving-and-parameter-management:4}
D'un seul coup, dix stratégies avaient l'erreur, et le premier des dix moniteurs s'est déclenché après cinq minutes~; avec un déploiement progressif, seul le moniteur du canari pouvait se déclencher, et sa session calme a duré quinze.
Mais d'un seul coup, dix stratégies ont accumulé de l'exposition jusqu'à la détection, contre une seule en déploiement progressif~: 1.35~million
de dollars-minute contre 53\,000.
\end{solution}

\begin{solution}{exo:pl:signal-serving-and-parameter-management:5}
Ce que la stratégie a utilisé est la valeur telle qu'elle était connue alors (75 points de base)~: c'est ce qui explique son comportement. Ce qu'elle aurait dû utiliser est la valeur
telle qu'elle est connue maintenant (60)~: c'est ce dont ont besoin la correction et tout retraitement du risque ou du P\&L.
\end{solution}

\begin{solution}{exo:pl:signal-serving-and-parameter-management:6}
Un contrôle de dérive comparant ce que chaque serveur exécute à ce qui est déclaré aurait révélé l'ancien code du huitième serveur~; retirer les \omterm{def:pl:signal-serving-and-parameter-management:flag}{indicateurs de fonctionnalité} plutôt que de les réutiliser,
et un déploiement progressif suivi serveur par serveur, auraient limité les dégâts. Un schéma sur les valeurs de paramètres n'aurait pas suffi~: la valeur
de l'\omterm{def:pl:signal-serving-and-parameter-management:flag}{indicateur de fonctionnalité} était valide, mais son sens sur ce serveur ne l'était pas.
\end{solution}

\begin{solution}{exo:pl:signal-serving-and-parameter-management:7}
Avec q10 comme canari, son moniteur ne se déclenche pas dans les trente minutes, les neuf autres reçoivent la valeur à quarante minutes, et le premier de
leurs moniteurs se déclenche cinq minutes plus tard~: détection 35 minutes après le changement, un pic de 1.31~million de dollars et 5.44~millions de dollars-minute --- pire
que d'un seul coup. Un canari dont le moniteur ne voit pas l'erreur ne fait que retarder le déploiement~; le canari doit être une stratégie sur laquelle l'erreur
apparaîtrait.
\end{solution}

\begin{solution}{exo:pl:signal-serving-and-parameter-management:8}
Versionné et relu, oui~; typé et contrôlé, non~: un relecteur voit «~50~» dans un diff et rien n'indique dans quelle unité, s'il est dans les bornes ou
de combien il a bougé~; le fichier dit ce qui est déclaré, pas ce que chaque stratégie exécute~; et une fusion part partout d'un coup. Un schéma, des approbations
proportionnelles à la taille, un contrôle de dérive et un déploiement fondé sur des \omterm{def:pl:signal-serving-and-parameter-management:flag}{indicateurs de fonctionnalité} restent nécessaires.
\end{solution}

\begin{solution}{pb:pl:signal-serving-and-parameter-management:1}
\begin{enumerate}
\item Une unité canonique (une fraction du capital), des unités acceptées (points de base, pourcentage, fraction), des bornes de 0 à 2~\%, un pas maximal de 50
points de base, et deux approbateurs au-delà de 20.
\item Refusée (aucune unité)~; refusée (hors des bornes et avec un pas trop grand)~; acceptée comme 0.005.
\item Toute personne autre que l'auteur~; un seul, ou deux au-dessus du seuil, distincts.
\item Temps de validité et temps de connaissance~: ce qui était en vigueur à un instant, tel qu'il était connu à un autre instant.
\item Que toute modification d'une ligne enregistrée soit détectée, et où.
\item La valeur de repli déclarée, avec le statut «~stale-fallback~» (ou «~missing-fallback~» s'il n'a jamais été publié).
\item Le responsable du signal, lorsque le signal est conçu.
\item Déployer un changement sur des stratégies choisies, puis revenir en arrière~; une source de dérive et de comportement latent.
\item Une stratégie qui exécute 0.5 là où le \omterm{def:pl:signal-serving-and-parameter-management:store}{magasin de paramètres} déclare 0.005.
\item Des ordres portant un \omterm{def:pl:signal-serving-and-parameter-management:flag}{indicateur de fonctionnalité} réutilisé pour du nouveau code ont déclenché l'ancien code sur un serveur que le nouveau code n'avait pas atteint.
\item Deux lots et deux cents lots.
\item Il se déclenche dans la première fenêtre de cinq minutes après le changement dont les exécutions dépassent 1.3 fois le nombre habituel~; il connaît le compte contrefactuel
de chaque stratégie, ce qu'un moniteur réel ne fait pas.
\item Après 5 minutes pour cinq stratégies, 15 pour trois, 30 pour une et 50 pour une.
\item Cinq minutes~; 410\,000~\$ au pic et 1.35~million de dollars-minute.
\item Quinze minutes~; 30\,000~\$ et 53\,000 dollars-minute.
\item \textbf{Résultat nommé.} D'un seul coup~: détection après 5 minutes, avec 410\,000~\$ d'inventaire excédentaire au pic et 1.35~million
de dollars-minute~; déploiement progressif~: 15 minutes, 30\,000~\$ et 53\,000 dollars-minute~; contrôle du schéma~: refusé à l'entrée, aucune exposition.
\item La détection ne compte qu'après les dégâts~; ce qui limite les dégâts, c'est le nombre de stratégies qui conservent l'erreur jusque-là, et si elle peut entrer
du tout.
\item Le schéma avec les unités~: c'est le moins coûteux, et il bloque toute une classe d'erreurs.
\item Des bornes plus serrées par stratégie, un pas maximal par rapport à la valeur courante, et une comparaison de l'effet de la nouvelle valeur dans un simulateur
avant le déploiement.
\item Une unité, des bornes, un responsable, des approbations, un historique et un plan de déploiement.
\end{enumerate}
\end{solution}

\begin{solution}{iq:pl:signal-serving-and-parameter-management:1}
Rendre l'unité partie intégrante de la valeur~: des paramètres typés avec une unité canonique, une conversion à l'entrée, des bornes et des limites de pas dans un schéma, et le refus
de toute valeur sans unité.
\iqlookfor{Types et schémas, plutôt que de la vigilance.}
\end{solution}

\begin{solution}{iq:pl:signal-serving-and-parameter-management:2}
Le signal devient périmé~; le \omterm{def:pl:signal-serving-and-parameter-management:service}{service de signaux} sert la valeur de repli déclarée, avec un statut sur lequel la stratégie peut agir, et une alerte est envoyée au responsable du signal.
La valeur de repli et la borne de fraîcheur font partie du contrat du signal.
\iqlookfor{Bornes de fraîcheur et valeurs de repli déclarées.}
\end{solution}

\begin{solution}{iq:pl:signal-serving-and-parameter-management:3}
Par le \omterm{def:pl:signal-serving-and-parameter-management:store}{magasin de paramètres}, avec deux approbations, derrière un \omterm{def:pl:signal-serving-and-parameter-management:flag}{indicateur de fonctionnalité}~: sur une stratégie canari où l'effet du changement est visible, surveillée par rapport
à son comportement habituel, puis sur quelques-unes, puis sur toutes, avec retour arrière automatique dès qu'un moniteur se déclenche.
\iqlookfor{Déploiement progressif avec un canari pertinent et retour arrière.}
\end{solution}

\begin{solution}{iq:pl:signal-serving-and-parameter-management:4}
L'écart entre la configuration déclarée et la configuration en cours d'exécution~; le détecter en faisant remonter à chaque processus la configuration qu'il a chargée (avec un
hash) et en la comparant au magasin, en continu et au démarrage.
\iqlookfor{Configuration en cours d'exécution remontée et comparée.}
\end{solution}

\begin{solution}{iq:pl:signal-serving-and-parameter-management:5}
Un magasin de valeurs typées, contrôlées par schéma, bitemporelles, avec des responsables~; des propositions et des approbations selon la taille du changement~; une piste d'audit à chaînage de hash~; des \omterm{def:pl:signal-serving-and-parameter-management:flag}{indicateurs de fonctionnalité} pour
des déploiements progressifs avec retour arrière~; une détection de dérive par rapport aux processus en cours d'exécution~; un \omterm{def:pl:signal-serving-and-parameter-management:service}{service de signaux} avec versions, fraîcheur et valeurs de repli~; et
des requêtes post-mortem sur ce qui était en production, et à quel moment.
\iqlookfor{Types, approbations, historique, déploiement, dérive.}
\end{solution}
